\section{Methodology}
Two classification baselines are evaluated: Logistic Regression and
XGBoost. Logistic Regression provides a linear baseline, while
XGBoost models nonlinear relationships and feature interactions.

The pipeline uses chronological training, validation, and test
partitions. Model selection and operating-threshold selection are
performed without using the test partition as a tuning set.
The saved evaluation identifies sigmoid calibration for both models.

Performance is measured using Precision--Recall Area Under the Curve
(PR-AUC) and Receiver Operating Characteristic Area Under the Curve
(ROC-AUC). Probability quality is assessed using the Brier score and
Expected Calibration Error (ECE). Lower Brier and ECE values indicate
better probabilistic accuracy and calibration under their respective
definitions.

The operating policy targets at least 60\% precision, with the
threshold selected on validation data. If no validation threshold
meets this target, threshold-dependent test metrics are unavailable;
they must not be reported as zero. For Logistic Regression, test
precision and recall are computed independently at the selected
validation threshold.
